\section{Experimental Evaluation}
\label{sec:evaluation}

We evaluate MA-ZeroPhish across five
dimensions: zero-day detection,
adaptive investigation, provenance-aware
reconciliation, selective collaboration,
and independent adjudication. The
experiments assess detection accuracy,
evidence robustness, and investigation
cost under unseen campaigns and
incomplete evidence.

\subsection{Experimental Setup}
\label{sec:exp-setup}

\noindent\textbf{Dataset and Zero-Day
Evaluation Protocol.}
The dataset comprises benign and phishing
URLs, webpages, SMS messages, and emails.
Chronological, campaign-disjoint splits
separate development, calibration, and
testing. Known campaign variants and
related domains or infrastructure are
grouped to minimize cross-split leakage.
Test campaigns are excluded from
few-shot examples, trigger configuration,
and estimator calibration.

Reputation and blacklist membership
are excluded as detection features.
We additionally evaluate samples
confirmed absent from selected reputation
sources at their observation time.
Where model cutoff dates are available,
we report a post-cutoff subset.
These controls strengthen zero-day
evaluation without claiming complete
exclusion from model pretraining.

\medskip
\noindent\textbf{Implementation and
Configuration.}
MA-ZeroPhish implements Phases 1--4 of Section~\ref{sec:process} using fixed model versions, declared
tools, and a common evidence schema.
The case-wide budget, dispatch costs,
trigger weights, confidence-band rules,
collaboration threshold $\tau$, round
limit $r_{\max}^{\mathrm{coll}}$, and
per-round reference limit $k$ are
fixed before testing.

The stopping-error estimator is trained
and calibrated on held-out intermediate
states assessed by a frozen Judge
against independently established
ground truth. Abstentions are evaluated
separately from substantive classification
errors. All methods use identical test
samples and acquisition conditions
unless explicitly varied.

\medskip
\noindent\textbf{Baselines.}
We compare MA-ZeroPhish with a
single-agent multimodal detector,
fixed all-specialist execution,
full multi-agent debate, and fixed
multi-agent execution without revision.
Baselines receive equivalent applicable
evidence, with tool access, model calls,
and resource budgets explicitly reported.
Both matched-budget and unrestricted-cost
comparisons distinguish architectural
benefits from additional computation.

\medskip
\noindent\textbf{Evaluation Metrics.}
Detection is assessed using precision,
recall, F1-score, false-positive rate,
and area under the precision--recall
curve. We report substantive-verdict
coverage, insufficient-evidence rate,
and selective risk to account for
abstention. Investigation efficiency
is measured by latency, model calls,
input/output tokens, acquisition
requests, and monetary cost.

Results include repeated-run variability
and paired confidence intervals on
identical test samples. Calibration
is assessed using reliability plots
and Brier score on intermediate states
receiving substantive Judge verdicts,
with abstention coverage reported
separately.


\subsection{Experiment 1: Zero-Day
Detection Performance}
\label{sec:exp1}

This experiment evaluates detection
performance on temporally separated,
campaign-disjoint phishing samples
and benign controls. MA-ZeroPhish
is compared with the single-agent,
fixed multi-agent, and full-debate
baselines using the same test set.
Results are reported separately
for URLs, webpages, SMS, and emails,
including message submissions
containing multiple URLs.

We report precision, recall,
F1-score, false-positive rate,
and precision--recall curves.
Because MA-ZeroPhish supports an
insufficient-evidence outcome,
selective risk and coverage are
reported alongside conventional
classification metrics. An additional
analysis isolates samples lacking
known reputation entries at the
time of observation.

This experiment determines whether
adaptive investigation and independent
adjudication maintain detection
performance on previously unseen
campaigns without relying on
reputation-based shortcuts.

\subsection{Experiment 2: Adaptive
Acquisition and Specialist Selection}
\label{sec:exp2}

This experiment evaluates Phase 1
under varying evidence availability
and acquisition budgets. We compare
adaptive specialist dispatch with
a fixed workflow that executes
all applicable specialists. Test
conditions include complete evidence,
partial browser captures, unavailable
network metadata, and message
submissions with multiple URLs.

The acquisition budget and number
of applicable modalities are varied
while recording selected specialists,
trigger coverage, acquisition success,
model calls, latency, and total cost.
We also measure the proportion of
initially unselected specialists
subsequently dispatched by the
Moderator.

The analysis examines the trade-off
between reduced initial investigation
cost and additional collaboration
required to recover omitted
modalities or missing evidence.

\subsection{Experiment 3: Provenance-Aware
Evidence Reconciliation}
\label{sec:exp3}

This experiment isolates the effect
of Phase 3 common-cause reconciliation.
We construct evaluation cases containing
repeated artifact citations, overlapping
cross-modal observations, independently
acquired corroboration, and conflicting
specialist interpretations. Dependency
ground truth is established from
acquisition records and manually
verified causal relationships where
available.

We compare full provenance-aware
reconciliation with variants that
(i) treat all agreeing observations
as independent and (ii) group
observations solely by semantic
similarity. Dependency identification
precision and recall, duplicate-support
rate, escalation frequency, and
downstream decision changes are
reported.

The experiment tests whether the
provenance graph discounts demonstrably
dependent evidence without incorrectly
merging independent observations
that express similar conclusions.

\subsection{Experiment 4: Uncertainty-Driven
Collaboration}
\label{sec:exp4}

This experiment evaluates the calibrated
stopping-error gate and targeted
collaboration in Phase 3. MA-ZeroPhish
is compared with no collaboration,
fixed-round collaboration, and
full multi-agent debate. We vary
the escalation threshold $\tau$,
round limit, and per-round citation
limit $k$ while maintaining the
same initial specialist records.

We measure detection performance,
selective risk, collaboration rounds,
specialist reinvocations, model calls,
tokens, latency, and total cost.
The stopping-error estimator is
evaluated using calibration plots,
Brier score, and discrimination
metrics on held-out intermediate
investigation states.

We additionally report escalation
precision: the proportion of
initiated rounds that resolve an
identified issue or yield new
eligible evidence. The results
characterize the accuracy--cost
trade-off and determine whether
the calibrated gate avoids
unnecessary full-agent debate.

\subsection{Experiment 5: Robustness
to Missing and Conflicting Evidence}
\label{sec:exp5}

This experiment evaluates the complete
framework under incomplete or
contradictory evidence. Controlled
test conditions progressively remove
served HTML, rendered DOM, screenshots,
or network metadata, and introduce
conflicting observations from
different modalities. Failed
acquisition attempts and stale
captures are recorded explicitly
rather than represented as benign
evidence.

We measure substantive-verdict
coverage, selective risk,
insufficient-evidence rate,
false-positive rate, recovery
through authorized local acquisition,
and the proportion of unresolved
issues disclosed in final decisions.
We also inspect whether the Judge
cites only eligible observations
and preserves material coverage
limitations.

This experiment examines whether
MA-ZeroPhish can abstain when
evidence is inadequate rather than
forcing agreement among specialists
or converting acquisition failures
into benign indicators.

\subsection{Experiment 6: Ablation Study}
\label{sec:exp6}

To quantify the contribution of
each mechanism, we evaluate the
full framework against five
ablated configurations:
\begin{enumerate}
\item Without adaptive specialist
selection, executing all applicable
agents initially.
\item Without common-cause
reconciliation, treating validated
observations as independent.
\item Without calibrated escalation,
using a fixed collaboration policy.
\item Without targeted collaboration,
replacing selective requests with
full-agent debate.
\item Without independent
adjudication, allowing the final
decision maker to observe specialist
verdicts and confidence bands.
\end{enumerate}

All configurations use identical
evaluation samples and fixed model
versions. Detection metrics,
selective risk, model calls,
tokens, latency, and cost are
reported for each ablation.
Additional analysis measures
revision-validation failures,
unsupported evidence citations,
and changes in final decisions
caused by removing the corresponding
mechanism.

The ablation study distinguishes
the benefits of adaptive routing,
provenance handling, calibrated
collaboration, and judgment
independence, rather than attributing
all improvements to multi-agent
execution alone.



% ---- properties --------------------------------------------------------

